% TC35 — handoff ข้ามบทบาทหลัง hardening
\begin{testscenario}{35}{4}
  \tcmeta{TC35}{CS 2568/01}{ระบบเว็บแอปพลิเคชัน ``NUCARE''}{นายสมชาย ใจดี}{วันที่ 7 พฤศจิกายน พ.ศ.\ 2569}{1.0}
  \tcname{การเชื่อมโยงบทบาทหลังแก้จุดขาด handoff}
  \tcpre{บัญชีทดสอบครบตาม \tccode{data.md} · unit ที่เกี่ยวข้องผ่าน}
  \tcdesc{ตรวจว่า SOS / คัดกรอง HIGH / สัญญาณแชทวิกฤต / แจ้งเตือนนิสิต / referral คนละ target ทำงานเชื่อมกัน}

  \tccase{1}{SOS มีผู้รับ}{%
    \begin{tcsteps}
      \item นิสิตกด SOS · MH รับ/ปิดที่หน้าหลัก
    \end{tcsteps}}{%
    มีคิวช่วยเหลือด่วนและแจ้งที่ปรึกษา}{%
    TC33 ผ่าน}{Pass}

  \tccase{2}{HIGH สร้างคิว counseling}{%
    \begin{tcsteps}
      \item ผลแบบประเมิน HIGH
    \end{tcsteps}}{%
    MH มีคิวรอรับจากระบบ}{%
    TC34 ผ่าน}{Pass}

  \tccase{3}{นิสิตเห็นการแจ้งเตือน}{%
    \begin{tcsteps}
      \item SD ส่ง reminder · เปิดหน้าหลักนิสิต
    \end{tcsteps}}{%
    มีกล่องการแจ้งเตือนและ toast ไม่กลืนรายการเดิม}{%
    inbox + DAI unit ผ่าน}{Pass}

  \tccase{4}{ส่งต่อคนละ target พร้อมกันได้}{%
    \begin{tcsteps}
      \item ที่ปรึกษาส่งสวัสดิการแล้วยังส่งศูนย์ให้คำปรึกษาได้
    \end{tcsteps}}{%
    ไม่ถูกบล็อกข้าม target}{%
    AR-T01/02 ผ่าน}{Pass}
\end{testscenario}
